\chapter{Excel Configuration}
\label{ch:excel}

All model inputs are defined in a single Excel workbook (\texttt{.xlsm} or \texttt{.xlsx}). This chapter describes each sheet and its required columns.

\section{Sheet Overview}

\begin{table}[h]
\centering
\begin{tabular}{lll}
\toprule
\textbf{Sheet Name} & \textbf{Purpose} & \textbf{Required} \\
\midrule
\texttt{Configuration} & Time range, elements, flags & Yes \\
\texttt{1\_1\_Definition\_Flows} & Flow definitions and connections & Yes \\
\texttt{1\_2\_Data\_Flows} & Flow time series data & Yes \\
\texttt{2\_1\_Definition\_Processes} & Process definitions and logic & Yes \\
\texttt{2\_2\_static\_TCs} & Static transfer coefficients & Optional \\
\texttt{2\_3\_dynamic\_TCs} & Time-varying transfer coefficients & Optional \\
\texttt{2\_4\_Initial\_Stock} & Initial stock values & Optional \\
\texttt{3\_1\_Definition\_DSM} & DSM parameters & If DSM used \\
\texttt{3\_2\_Definition\_FOMP} & FOMP parameters & If FOMP used \\
\texttt{4\_1\_Uncertainty\_Parameters} & Monte Carlo distributions & If MC used \\
\texttt{5\_1\_Scenario\_Definitions} & Scenario configurations & If scenarios used \\
\bottomrule
\end{tabular}
\caption{Excel workbook sheets.}
\label{tab:sheets}
\end{table}

\section{Configuration Sheet}

Defines global model settings.

\begin{table}[h]
\centering
\begin{tabular}{llp{6cm}}
\toprule
\textbf{Parameter} & \textbf{Example} & \textbf{Description} \\
\midrule
\texttt{Start\_Year} & 2020 & First year of the model \\
\texttt{End\_Year} & 2050 & Last year of the model \\
\texttt{Elements} & material,WC,DM,CC & Comma-separated element list \\
\texttt{Element\_ID\_X} & CC & Element name for hierarchy \\
\texttt{Parent\_Element\_ID\_X} & DM & Parent element name \\
\bottomrule
\end{tabular}
\caption{Configuration sheet parameters.}
\end{table}

\section{Flow Definitions (Sheet 1\_1)}

Defines all flows in the system: their source, destination, and element composition.

\subsection{Required Columns}

\begin{table}[h]
\centering
\begin{tabular}{llp{6cm}}
\toprule
\textbf{Column} & \textbf{Example} & \textbf{Description} \\
\midrule
\texttt{Flow\_ID} & \texttt{F\_01\_02} & Unique flow identifier \\
\texttt{Source\_Process} & 1 & ID of the starting process \\
\texttt{Target\_Process} & 2 & ID of the ending process \\
\texttt{E1\_name} & WC{[}\%{]} & First element name (composition fraction) \\
\texttt{E1\_value} & 60 & Fraction of material that is this element (\%) \\
\texttt{E2\_name} & CC\_DM{[}\%{]} & Second element (CC as \% of DM) \\
\texttt{E2\_value} & 45 & Fraction value \\
\bottomrule
\end{tabular}
\caption{Flow definition columns.}
\end{table}

\subsection{Naming Convention}

Flow IDs follow the pattern \texttt{F\_<source>\_<target>\_<name>}:
\begin{itemize}
  \item \texttt{F\_01\_02} --- Flow from Process 1 to Process 2
  \item \texttt{F\_05\_00\_CO2} --- CO\textsubscript{2} flow from Process 5 to the environment
\end{itemize}

\section{Flow Data (Sheet 1\_2)}

Contains time series values for input flows (flows from the system boundary).

\begin{table}[h]
\centering
\begin{tabular}{llllll}
\toprule
\textbf{Flow\_ID} & \textbf{Year} & \textbf{E1\_value} & \textbf{Unit} \\
\midrule
\texttt{F\_00\_01} & 2020 & 1000 & Mg \\
\texttt{F\_00\_01} & 2030 & 1500 & Mg \\
\texttt{F\_00\_01} & 2050 & 2000 & Mg \\
\bottomrule
\end{tabular}
\caption{Flow data example. Years between data points are interpolated linearly.}
\end{table}

\begin{tipbox}
You only need to provide data for key years. BioDYM interpolates linearly between them. For example, defining values at 2020, 2030, and 2050 will automatically fill all intermediate years.
\end{tipbox}

\section{Process Definitions (Sheet 2\_1)}

Defines each process node in the system.

\begin{table}[h]
\centering
\begin{tabular}{llp{5.5cm}}
\toprule
\textbf{Column} & \textbf{Example} & \textbf{Description} \\
\midrule
\texttt{ID} & 3 & Unique process number \\
\texttt{Name} & Use\_Phase & Human-readable name \\
\texttt{Process\_Logic} & Splitter & How the process handles material \\
\texttt{Stock\_Configuration} & Stock & Whether the process has a stock \\
\bottomrule
\end{tabular}
\caption{Process definition columns.}
\end{table}

\subsection{Process Logic Values}

\begin{itemize}
  \item \textbf{Input} --- System boundary (Process 0); source of input flows
  \item \textbf{Splitter} --- Distributes inflow to outflows by TCs; preserves composition
  \item \textbf{Transformer} --- Like Splitter, but allows element-specific TCs (changes composition)
  \item \textbf{Pass-through} --- Passes all material directly to a single outflow
  \item \textbf{DSM} --- Dynamic stock model with lifetime distributions
  \item \textbf{FOMP} --- First-order multi-pool organic matter decay
\end{itemize}

\subsection{Stock Configuration Values}

\begin{itemize}
  \item \textbf{NoStock} --- Process has no stock
  \item \textbf{Stock} --- Process accumulates stock (zero initial stock)
  \item \textbf{Stock\_with\_InitialStock\_Decay} --- Stock with initial value, simple exponential decay
  \item \textbf{Stock\_with\_InitialStock\_Cohort} --- Stock with initial value, ODYM age-cohort method
\end{itemize}

\section{Transfer Coefficients}

TCs control how material leaving a process is distributed among its outflows.

\subsection{Static TCs (Sheet 2\_2)}

Constant values that don't change over time:

\begin{table}[h]
\centering
\begin{tabular}{llll}
\toprule
\textbf{TC\_ID} & \textbf{Value} & \textbf{Unit} \\
\midrule
\texttt{TC\_05\_06} & 0.7 & 1 \\
\texttt{TC\_05\_07} & 0.3 & 1 \\
\bottomrule
\end{tabular}
\caption{Static TC example. Outflows from Process 5: 70\% to Process 6, 30\% to Process 7.}
\end{table}

\subsection{Dynamic TCs (Sheet 2\_3)}

Time-varying values, defined at key years and interpolated:

\begin{table}[h]
\centering
\begin{tabular}{llll}
\toprule
\textbf{TC\_ID} & \textbf{Year} & \textbf{Value} \\
\midrule
\texttt{TC\_05\_06} & 2020 & 0.6 \\
\texttt{TC\_05\_06} & 2040 & 0.8 \\
\texttt{TC\_05\_07} & 2020 & 0.4 \\
\texttt{TC\_05\_07} & 2040 & 0.2 \\
\bottomrule
\end{tabular}
\caption{Dynamic TC example. Values interpolated between 2020 and 2040.}
\end{table}

\subsection{TC Naming Convention}

\begin{itemize}
  \item \textbf{Standard}: \texttt{TC\_<source>\_<target>} --- e.g., \texttt{TC\_05\_06}
  \item \textbf{Element-specific}: \texttt{TC\_E<idx>\_<source>\_<target>} --- e.g., \texttt{TC\_E2\_05\_06} for the WC-specific TC
\end{itemize}

Element-specific TCs are only used with \emph{Transformer} processes. Splitters use the same TC for all elements.

\begin{warnbox}
TCs for all outflows of a process must sum to 1.0 (100\%) at every time step. BioDYM normalizes automatically if they don't, but will print a warning if the deviation exceeds 5\%.
\end{warnbox}

\section{Initial Stock (Sheet 2\_4)}

Defines starting stock values for processes with stock configuration:

\begin{table}[h]
\centering
\begin{tabular}{llll}
\toprule
\textbf{Process\_ID} & \textbf{Parameter\_Name} & \textbf{Parameter\_Value} & \textbf{Unit} \\
\midrule
1 & Initial\_Stock\_material & 5000 & Mg \\
1 & Initial\_Stock\_WC{[}\%{]} & 15 & \% \\
1 & Initial\_Stock\_DM{[}\%{]} & 85 & \% \\
1 & Initial\_Stock\_CC{[}\%{]} & 45 & \% \\
1 & Annual\_Consumption\_Rate & 0.05 & 1/year \\
\bottomrule
\end{tabular}
\caption{Initial stock configuration example.}
\end{table}
